% ============================================
% Discussion
% ============================================
\section{Discussion}
\label{sec:discussion}

Our work demonstrates that interspecies interaction strength governs community coalescence outcomes. In both experiments and modeling, we showed that in the presence of intermediate or strong interspecies interactions, community-level selection (CLS) was the most common outcome, indicating that the community becomes the unit of selection. In contrast, when interactions were weak, Mixture dominated, species behaved more independently, and CLS disappeared. Theoretical modeling with a gLV framework predicts a shift from CLS to Mixture as interspecies interaction strength decreases. Varying nutrient concentration provided an experimental control for interaction strength, shifting outcome patterns in a predictable manner consistent with the model. Finally, species-level predictability revealed two distinct regimes of community-level selection: a species-driven regime, where dominant pH-modifying taxa largely determine the winning community, and an emergent regime, where multi-species collective dynamics reduce predictability. Together, our work highlights the versatile and variable nature of community coalescence, while community-level selection and its predictability emerge conditionally.

Our work provides the first experimental demonstration of alternative regimes in community coalescence, reconciling previously contrasting observations of community-level versus species-level selection. Prior studies have reported both outcomes: some found communities behaving as cohesive units~\citep{Tikhonov2016, DiazColunga2022}, while others observed species responding independently~\citep{VanderGucht2007, Vermeij1991}. Our results suggest that these conflicting observations may reflect differences in interaction strength across experimental systems. For instance, the absence of community-level selection in \textit{in vitro} gut microbiome coalescence experiments~\citep{Goldman2025, Walton2025} is consistent with our framework. Rich media such as BHI are analogous to our Nutr$-$ condition, providing complex and diverse resources while lacking high concentrations of readily metabolizable carbon sources. This resource structure reduces both direct resource competition and pH-mediated interactions, thereby weakening interspecies interactions~\citep{Kehe2021} and favoring species-level rather than community-level dynamics. Likewise, interaction strength may vary systematically across ecosystems: microbial biofilms experience strong metabolic interactions due to high cell densities and shared resources~\citep{Nadell2016}, whereas mobile macroscopic organisms interact more transiently across larger spatial scales. These fundamental differences between microbial and macroscopic systems may explain why community-level selection is more frequently observed in microbial coalescence studies. This framework provides a unified perspective: rather than asking whether communities are cohesive units, we should ask under what conditions they become so, and interaction strength emerges as a key determinant.


A growing body of work has identified interaction strength as a key control parameter governing diverse aspects of microbial community dynamics, including diversity \citep{Hu2022, Marsland2019}, stability \citep{Ratzke2020}, and invasion outcomes \citep{Hu2025, Kurkjian2021}. Our results extend this claim to include coalescence, providing further evidence that interaction strength determines when communities exhibit collective behavior. \jy{The parallel with priority effects is particularly instructive; recent work showed that assembly order shapes final composition only under strong interactions, whereas weak interactions lead to convergent outcomes regardless of assembly history \citep{Hu2025, Fukami2015}.} Our coalescence results follow an analogous pattern---strong interactions preserve parental identity, while weak interactions yield mixtures independent of community origin. Together, these findings reinforce interaction strength as a coarse-grained parameter that predicts history-dependent community dynamics across multiple ecological contexts.

In nature, coalescence occurs in richer settings that may shift regimes and outcomes. Environmental heterogeneity (temperature, pH buffering, and other physicochemical factors) co-varies with interaction strength and can reshape coalescence across habitats \citep{Rillig2015, Castledine2020, Rillig2017, Tropini2017}. In host-associated microbiomes, host filtering, priority effects, and spatial structure further influence successful colonization, making coalescence a useful framework for predicting and steering microbiome transfer \citep{Liu2024, Smillie2018, Amor2024}. Our experiments focus on steady-state outcomes and thus do not capture temporal dynamics such as fluctuating resources, migration pulses, or host responses. Additionally, both our experimental system and theoretical model are dominated by competitive interactions; systems with substantial mutualism or facilitation may exhibit qualitatively different dynamics. Incorporating these axes into both theory and experiment will help build a more comprehensive picture of community-level selection across habitats and scales.

